\subsection{光滑代数}

引理 5.—— 设 \(\rho: A \to B\) 为 Noether 局部环之间的局部同态。假设 B 在 A 上是本质有限型的。为使 A-代数 B 形式光滑，必须且只须 A-模 B 平坦，且 \(\kappa_{A}\)-代数 \(\kappa_{A} \otimes_{A} B\) 绝对正则。

存在整数 \(n \geqslant 0\)、\(A[T_{1}, \ldots, T_{n}]\) 的素理想 q，以及从 \(A[T_{1}, \ldots, T_{n}]_{q}\) 到 B 的满射同态 h。用 C 表示局部 A-代数 \(A[T_{1}, \ldots, T_{n}]_{q}\)；它是形式光滑的（\(n^{\circ} 3\)，例 2 及 \(n^{\circ} 2\)，命题 4, a)）且在 A 上平坦，并且可把 B 等同于 A-代数 C/J，其中 \(J = \text{Ker}(h)\)。

令 \(\overline{\mathrm{C}} = \kappa_{\mathrm{A}} \otimes_{\mathrm{A}} \mathrm{C}\)，\(\overline{\mathrm{B}} = \kappa_{\mathrm{A}} \otimes_{\mathrm{A}} \mathrm{B}\)。假设 B 在 A 上形式光滑。于是 \(\kappa_{\mathrm{A}}\)-代数 \(\overline{\mathrm{B}}\) 形式光滑（第 2 节，命题 4, b)），故绝对正则（第 5 节，定理 2 的推论 2）。此外，由于 \(\overline{\mathrm{C}} / \overline{\mathrm{JC}}\) 与 \(\overline{\mathrm{B}}\) 等同，且 \(\kappa_{\mathrm{A}}\)-代数 \(\overline{\mathrm{C}}\) 形式光滑，故理想 \(\overline{\mathrm{JC}}\)（在 \(\overline{\mathrm{C}}\) 中）是完全割的（第 5 节，定理 2）。于是由 § 5, 第 6 节，命题 6 可知 A-模 B 平坦。

反之，假设 B 在 A 上平坦，且 \(\kappa_{\mathrm{A}}\)-代数 \(\overline{\mathbf{B}}\) 绝对正则。于是 \(\kappa_{\mathrm{A}}\)-局部代数 \(\overline{\mathbf{B}}\) 形式光滑（第 9 节的注 1，取 \(k = k_0 = \kappa_{\mathrm{A}}\)）。令 \(\overline{\mathbf{J}} = \kappa_{\mathrm{A}} \otimes_{\mathrm{A}} \mathrm{J}\)；由于 B 是平坦 A-模，典范映射 \(\overline{\mathbf{J}} \to \overline{\mathbf{J}}\overline{\mathbf{C}}\) 是双射，且 \(\overline{\mathbf{B}}\) 与 \(\overline{\mathbf{C}}/\overline{\mathbf{J}}\) 等同。由此（第 2 节的注 1）可知典范映射

\[
\overline {{\mathrm{J}}} / \overline {{\mathrm{J}}} ^ {2} \longrightarrow \overline {{\mathrm{B}}} \otimes_ {\overline {{\mathrm{C}}}} \Omega_ {\kappa_ {\mathrm{A}}} (\overline {{\mathrm{C}}})
\]

是单射且具有一个收缩。而 \(\overline{\mathrm{J}}/\overline{\mathrm{J}}^{2}\) 与 \(\kappa_{\mathrm{A}}\otimes_{\mathrm{A}}\mathrm{J}/\mathrm{J}^{2}\) 等同，故与 \(\overline{\mathrm{B}}\otimes_{\mathrm{B}}\mathrm{J}/\mathrm{J}^{2}\) 等同。另一方面，\(\overline{\mathrm{C}}\)-模 \(\Omega_{\kappa_{\mathrm{A}}}(\overline{\mathrm{C}})\) 典范地同构于 \(\overline{\mathrm{C}}\otimes_{\mathrm{C}}\Omega_{\mathrm{A}}(\mathrm{C})\)（A, III, 第 136 页，命题 20），故 \(\overline{\mathrm{B}}\otimes_{\overline{\mathrm{C}}}\Omega_{\kappa_{\mathrm{A}}}(\overline{\mathrm{C}})\) 典范地同构于 \(\overline{\mathrm{B}}\otimes_{\mathrm{C}}\Omega_{\mathrm{A}}(\mathrm{C})\)。对 \(\overline{\mathrm{B}}\) 的极大理想取商，即得一个单同态

\[
\kappa_ {B} \otimes_ {B} J / J ^ {2} \longrightarrow \kappa_ {B} \otimes_ {C} \Omega_ {A} (C)
\]

它正是 \(\bar{d}_{\kappa_{\mathrm{B}}}\)。于是 B 在 A 上形式光滑（定理 3）。

定理 4.—— 设 A 为 Noether 环，B 为本质有限型的 A-代数。下列条件等价：

\begin{enumerate}
\def\labelenumi{(\roman{enumi})}
\item
  A-代数 B 形式光滑；
\item
  对每个 \(\mathfrak{q} \in \operatorname{Spec}(\mathrm{B})\)，A-代数 \(\mathrm{B}_{\mathfrak{q}}\) 形式光滑（相应地，关于 \(\mathfrak{q}\mathrm{B}_{\mathfrak{q}}\)-进拓扑形式光滑）；
\item
  A-模 B 平坦，且对每个 \(\mathfrak{p} \in \operatorname{Spec}(A)\)，\(\kappa(\mathfrak{p})\)-代数 \(\kappa(\mathfrak{p}) \otimes_{A} B\) 绝对正则；
\item
  A-模 B 平坦，且对每个正则 A-代数 R，环 R ⊗\_A B 正则；
\item
  A-模 B 平坦，且同态 \(\mu:B\otimes_{A}B\to B\)（由 \(\mu(x\otimes y)=xy\) 定义）的核是完全割理想。
\end{enumerate}

(i) 与 (ii) 的等价性由第 9 节定理 3 的推论 2 得出。

(i)⇒(v)：假设 A-代数 B 形式光滑。设 q 为 B 的素理想，p 为它在 A 中的原像。\(A_{\mathfrak{p}}\)-代数 \(B_{\mathfrak{q}}\) 形式光滑（第 2 节的命题 4, a)），故平坦（引理 5）；因而 A-模 B 平坦（II, § 3, 第 4 节，命题 15）。另一方面，环 B ⊗A B 是 Noether 的（§ 6, 第 1 节，命题 2 的推论），故由第 5 节定理 2 的推论 3，理想 Ker μ 是完全割的。

(v)⇒(iii)：假设条件 (v) 满足。令 I = Ker(μ)。设 p ∈ Spec(A)。映射

\[
1 \otimes \mu : \kappa (\mathfrak {p}) \otimes_ {\mathrm{A}} (\mathrm{B} \otimes_ {\mathrm{A}} \mathrm{B}) \rightarrow \kappa (\mathfrak {p}) \otimes_ {\mathrm{A}} \mathrm{B}
\]

等同于映射

\[
\mu_ {\mathfrak {p}}: (\kappa (\mathfrak {p}) \otimes_ {\mathrm{A}} \mathrm{B}) \otimes_ {\kappa (\mathfrak {p})} (\kappa (\mathfrak {p}) \otimes_ {\mathrm{A}} \mathrm{B}) \rightarrow \kappa (\mathfrak {p}) \otimes_ {\mathrm{A}} \mathrm{B}
\]

它是由 \(\kappa(\mathfrak{p})\)-代数 \(\kappa(\mathfrak{p})\otimes_{\mathrm{A}}\mathrm{B}\) 的乘法导出的。理想 \(\operatorname{Ker}(\mu_{\mathfrak{p}})\) 与 \(\mathrm{I}(\kappa(\mathfrak{p})\otimes_{\mathrm{A}}(\mathrm{B}\otimes_{\mathrm{A}}\mathrm{B}))\) 等同。由于 A-模 B 平坦（\(\S 5\), 第 6 节，命题 6），它是完全割的。于是断言 (iii) 由 \(\S 6\) 第 5 节的命题 8 得出。

(iii)⇒(ii)：设 q 为 B 的素理想，p 为其在 A 中的原像。在 (iii) 的假设之下，A \(_{p}\)-模 B \(_{q}\) 平坦，而 κ(p)-代数 κ(p) ⊗A \(_{p}\) B \(_{q}\)（它与 κ(p) ⊗A B 的一个分式环等同）绝对正则（§ 6, 第 4 节，命题 6）。于是由引理 5，B \(_{q}\) 在 A \(_{p}\) 上形式光滑，故在 A 上形式光滑（第 2 节，命题 3 与 4）。

(iii)⇒(iv)：设在 (iii) 的假设之下。设 R 为正则 A-代数。R-模 R ⊗A B 平坦（I, § 2, 第 7 节，命题 8 的推论 2）。设 r 为 R 的素理想，p 为其在 A 中的原像；环 κ(r)⊗R(R⊗AB)（它与 κ(r)⊗κ(p) (κ(p)⊗A B) 等同）是正则的（§ 6, 第 4 节，命题 7 的推论 2）。于是环 R ⊗A B 正则（§ 4, 第 5 节，命题 9 的推论）。

(iv)⇒ (iii)：设 p 为 A 的素理想，k 为 κ(p) 的一个扩张；在 (iv) 的假设之下，环 \(k \otimes_{\kappa(\mathfrak{p})} (\kappa(\mathfrak{p}) \otimes_{\mathrm{A}} \mathrm{B})\)（它与 \(k \otimes_{A} B\) 等同）是正则的，由此得 (iii)。

定义 2.—— 设 A 为 Noether 环。若 A-代数 B 是本质有限型的且满足定理 4 的诸等价条件，则称它是光滑的。

命题 10.—— 设 A 为 Noether 环。

\begin{enumerate}
\def\labelenumi{\alph{enumi})}
\item
  设 A' 为 Noether A-代数，B 为光滑 A-代数。则 A'-代数 A' ⊗\_A B 光滑。
\item
  设 B 为光滑 A-代数，C 为光滑 B-代数。则 A-代数 C 光滑。
\item
  设 B 与 C 为两个光滑 A-代数。则 A-代数 B ⊗A C 光滑。
\end{enumerate}

这由第 2 节的命题 4 以及关于本质有限型代数的类似命题（§ 6, 第 1 节）得出。

例.—— 1) 域 \(k\) 上的光滑代数恰是本质有限型且绝对正则的 \(k\)-代数。

\begin{enumerate}
\def\labelenumi{\arabic{enumi})}
\setcounter{enumi}{1}
\tightlist
\item
  设 A 为 Noether 环，\(\mathbf{T} = (\mathrm{T}_i)_{i\in I}\) 为有限的未定元族。A-代数 A{[}T{]} 是光滑的。更一般地，设 \(\mathrm{F}_1,\ldots ,\mathrm{F}_m\) 为 A{[}T{]} 的元素，B 为 A-代数 A{[}T{]}/ \((\mathrm{F}_1,\dots ,\mathrm{F}_m)\)。若在 B 的每个极大理想 \(\mathfrak{n}\) 处，模 \(\mathfrak{n}\) 之下矩阵 \(\left(\frac{\partial\mathrm{F}_j}{\partial\mathrm{T}_i}\right)\) 的类都有秩 \(m\)，则 A-代数 B 光滑（第 9 节，定理 3）。
\end{enumerate}

